% Page 073

\seite{73}

einschließlich 7.\ Juli.
\margmark{13. Juli.}%
Infolge pflichtmäßigen Meldung wurden die Prüfungen am\ob
3.\ 4. u 5.\ Juli unterbrochen.

\darueber{gewählt}\ob
In Dasbach Heinrich der Wahl zum Gemeindevorsteher-Stell\ohb
vertreter, wurde an seiner statt Schirmers Nikolaus\ob
gewählt.

\margmark{24. September}%
.   In der Wahlvorbereitung zur vierjährigen Gemeinde\ohb
Drostwahl sowie die der Wahl selbst wurden Unregelmäßig\ohb
keiten vorgenommen und, die gegen die Gemeindevorschrift\ob
und die Wahlordnung verstoßen, müßte die dagegen seitens\ob
Einsprüche in aller erster Linie vom Kreisausschuß aner\ohb
kannt und die Wahl für ungültig erklärt werden. Darüber\ob
hire in absehbarer Zeit eine Nachricht der Gemeindevorsteher\ob
zu bringen. (Sieh Seite 77.)

\margmark{24. September}%
.   Die Großferien dauern vom 24.\ 9. bis einschließlich\ob
28.\ 10. (35 Tage.)

Schwere Großen- und Unwetter wurde diesen Sommer\ob
im Eifelkreise. Anfangs des Jahres die meiste Schneetreib auf\ob
Balken, hohen Halde, fing es an zu regnen. Ihre sogen.\ob
Wolkenbrüche (ungefähr 3 Stunde) fielen an den vorletzten\ob
großen Tal der Nims\unsicher Brücke\unsicher eine Flut, aufgehäufet\ob
wurde der Ballen bald mit Wasser abgerissen – die Be\ohb
richtigung wird – so bitten sie unter Pflanzungen. Aufwärts\ohb
lich wüteten die Garten und den Boden auseinanderge\ohb
zogen, einen großen Teil der Dörfer ganz abgeschwommen\ob
worden. Der Schaden ist groß – sehr groß. Die Dort\unsicher{} konnte\ob
werden.

\margmark{30. Oktober.}%
    Gegen der Anfechtung der Kreisausschüsse vgl. Ungültig\ohb
keitserklärung des vierjährigen Gemeinderatswahlresult wurde\ob
von den Gewählten beim Bezirksausschuß in Trier Berufung\ob
eingelegt.

\margmark{20. Dezember.}%
    Die Anfechtung der Kreisausschüsse vgl. Ungültigkeit\ob
der Gemeinderatswahlresulte wurde vom Bezirksausschuß\ob
anerkannt. Die Neuwahlung wurde vorgeschrieben. Dieselb\unsicher{} soll\ob
nun bald am Neujahr statt\unsicher.
\pend
